
\section{Design Spaces and Insights}
We conduct a comprehensive ablation study to provide more insights about our framework design, including token initialization, supervised training, reward function design, reinforcement learning training, sequence augmentation, and positional encoding.
The result is listed in Table~\ref{tab: ablation}.

\subsection{Training Recipe} \label{sec: ablation training recipe}

\textbf{Token Initialization.}
We explore three token initialization methods for \bk token, i.e., centroid (Cent.), context-based (Ctx.), and semantic anchor (Sem.) initializations. 
The detailed discussions about these methods are provided in Appendix~\ref{app: token initialization}.
According to Table~\ref{subtab: initialization}, both context and semantic-based strategies outperform the naive centroid method, with semantic initialization achieving the highest accuracy on both \textsc{Gsm8k} and \textsc{Math} datasets.
Notably, the performance gap between informed initialization methods, i.e., context and semantic, and the centroid method is more substantial on the challenging MATH benchmark, indicating that \textbf{semantically meaningful token embeddings provide greater benefits for complex reasoning tasks.}

\textbf{Supervised Fine-Tuning.} 
We compare mSFT against standard SFT using the augmented dataset $\gD_\text{aug}$. 
As shown in Table~\ref{subtab: sft}, standard SFT exhibits negative learning, evidenced by a performance decline compared to baseline (Few-Shot).
This degradation likely stems from error tokens within $\gD_\text{aug}$ leading the trained models to imitate error tokens and generate erroneous outputs. 
Conversely, mSFT mitigates this by masking error tokens, thereby yielding a performance gain.

 \begin{table*}[!t]
    \centering
    \caption{
        The performance comparison on benchmark math reasoning and coding datasets. The best performance is highlighted in bold, and the runner-up is underlined.
    } \label{tab: main}
    % \renewcommand{\arraystretch}{1}
    \setlength{\tabcolsep}{11pt}
    \resizebox{\linewidth}{!}{
    \begin{tabular}{ccccccccccccc}
    \toprule
    
    \multirow{2}{*}{\textbf{Method}} & \multicolumn{3}{c}{Llama-3.2-1B} & \multicolumn{3}{c}{Llama-3.1-8B} & \multicolumn{3}{c}{Qwen3-4B} \\
    \cmidrule(lr){2-4} \cmidrule(lr){5-7} \cmidrule(lr){8-10} 
    & \textsc{Gsm8k} & \textsc{Math} & \textsc{Mbpp} & \textsc{Gsm8k} & \textsc{Math} & \textsc{Mbpp} & \textsc{Gsm8k} & \textsc{Math} & \textsc{Mbpp} \\
    \midrule
    Few-Shot & 8.8 & 3.2 & 25.1 & 53.3 & 24.1 & 48.0 & 72.3 & 54.1 & 67.6 \\ 
    SFT & 14.3 & 10.5 & 29.9 & 55.6 & 28.1 & 58.8 & 81.5 & 60.7 & 70.2\\ 
    RFT & 26.1 & 12.7 & 29.5 & 59.1 & 30.9 & 58.6 & 80.3 & 63.8 & 71.3\\
    STaR+ & 26.5 & 12.9 & 30.7 & 59.3 & 30.6 & 59.1 & 82.1 & 64.0 & 70.9 \\ 
    ReVISE & \underline{28.1} & \underline{13.4} & \underline{33.3} & \underline{61.6} & \underline{33.6} & \underline{60.3} & \underline{84.2} & \underline{66.3} & \underline{71.9}\\ 
    \rowcolor{Gray} \textbf{Ours} & \textbf{31.3} & \textbf{15.2} & \textbf{34.1} & \textbf{63.4} & \textbf{36.2} & \textbf{61.1} & \textbf{86.1} & \textbf{69.5} & \textbf{72.5}\\
   \bottomrule 
    \end{tabular}
    }
\end{table*}


\textbf{Reward Function.} 
We examine reward function designs in Table~\ref{subtab: reward}. 
Both components contribute distinctly, i.e., $\mathcal{R}_\text{inc}$ provides a 10\% gain on \textsc{Gsm8k}, whereas $\mathcal{R}_\text{pen}$ yields a 12\% gain on \textsc{Math}. 
This suggests that \textbf{simpler tasks (\textsc{Gsm8k}) benefit primarily from outcome-based rewards}, where correct erasing naturally emerges from accuracy signals. 
Conversely, \textbf{complex reasoning tasks require explicit regularization of erase behavior to prevent degenerate strategies}, i.e., either over-reliance or under-utilization of the undo mechanism.

Besides our proposed reward functions, we also test the designs in BSAFE+~\cite{selreinforcement}, SCoRe~\cite{kumar2024training}, and MGRPO~\cite{ding2025multi}. 
Implementation details are provided in Appendix~\ref{app: additional reward}.
We find that our designed reward function performs better than those in SCoRe and MGRPO. 
Although BSAFE+ achieves comparable performance with our proposed method, it employs an external LLM as a critic to compute the reward, which introduces computational overhead. 
Therefore, our proposed method remains a better design in terms of performance and computational costs.



\textbf{GRPO.} 
We assess the efficacy of GRPO over base and fine-tuned models in Table~\ref{subtab: grpo}.
Note that, the SFT result is the GRPO training based on the model after SFT over $\gD_\text{aug}$.
GRPO consistently enhances performance across different settings, even for the SFT model that has shown negative learning in Table~\ref{subtab: sft}.
This indicates that \textbf{GRPO effectively contributes to \modelname's generate steps, leading to improved reasoning capabilities.}

\subsection{Augmentation}

\textbf{Sequence Augmentation.}
We examine the proposed sequence augmentation method in Section~\ref{sec: sequence augmentation}, with two additional sequence augmentation approaches for erase, i.e., stochastic error injection (Rand.) and hard sample sequence augmentation (Hard).  
Appendix~\ref{app: sequence augmentation} provides the details of these two augmentation methods.
Table~\ref{subtab: augmentation} demonstrates that our proposed method outperforms both alternatives, indicating that \textbf{targeted error injection focusing on critical reasoning steps is more effective than random or hard sample-based augmentations for training robust non-monotonic models.}

\textbf{Positional Encoding.}
As discussed in \citet{cundy2023sequencematch}, altering the positional encoding may also bring potential performance gains.
The intuition is that correction tokens should inherit the position indices of the error tokens they replace, such that the final semantic sequence appears positionally contiguous.
We investigate this approach in Table~\ref{subtab: position}, and we find that positional encodings consistently degrade performance.
We attribute this to RoPE~\citep{su2024roformer}, which rewinds position indices, introducing non-monotonic discontinuities that conflict with pretrained attention patterns.
Moreover, \textbf{standard monotonic positions already preserve the full context of the error-correction process}, allowing the model to learn correction semantics directly from the \bk token and its surrounding context rather than relying on positional cues.